\section{Introduction}
\label{sec:intro}

Molecular-dynamics studies of colloidal self-assembly readily generate more structural data
than can be assigned by hand. A single parameter sweep produces thousands of configurations,
many containing several coexisting phases, partially ordered regions, or structures that
were not anticipated when the study was designed. Assigning a structural label to every
domain is a prerequisite for mapping a phase diagram, and it does not scale manually.

Established approaches address parts of this problem. Bond-orientational order parameters
characterise local environments and underpin much of the crystal-detection literature
\TODO{VERIFY CITATION: Steinhardt et al.}, with averaged variants improving robustness to
thermal noise \TODO{VERIFY CITATION: Lechner and Dellago}. Supervised classifiers trained on
labelled environments perform well when the classes are known in advance
\TODO{VERIFY CITATION}. Unsupervised methods discover recurring local environments without
labels \TODO{VERIFY CITATION: Boattini et al.; Reinhart et al.}. Fixed descriptors such as
the smooth overlap of atomic positions provide a general-purpose representation of local
structure \TODO{VERIFY CITATION: Bartok et al.; DScribe}.

A workflow intended to operate on data containing genuinely unfamiliar structures needs
several properties that these components do not supply on their own. It must resolve a
heterogeneous frame into spatial domains before assigning structure, since a descriptor
averaged over two coexisting phases describes neither. It must distinguish domains that
resemble the training atlas from those that do not, rather than assigning every domain to
its nearest known class. When a domain is unfamiliar, it must be comparable against
physically interpretable references rather than only against the training data. And it must
be able to report that no reference is compatible, instead of forcing a choice.

We investigate such a workflow for binary colloidal crystals. The components are individually
standard; what we examine is whether they compose into a system that detects structures from
families absent at training time, identifies them against an independent reference library
built from crystallographic prototypes, and abstains when the correct reference is
unavailable. We report where this succeeds and where it does not: in particular, the
acceptance threshold is the least stable element of the detector, and the expansion stage of
the workflow is set up but not demonstrated.
